% Part: computability
% Chapter: computability-theory
% Section: rice-theorem

\documentclass[../../../include/open-logic-section]{subfiles}

\begin{document}

\olfileid{cmp}{thy}{rce}
\olsection{د رايس قضيه}

که لږ فکر وکړئ، وبه وينئ چې د~$\fn{Tot}$ ځانګړي جزئيات د
\olref[tot]{prop:total} په ثبوت کښې هېڅ رول نۀ لري۔ موږ~$h(x,y)$
داسې جوړ کړے و چې کټ مټ هغه وخت د ثابتې تابع~$j(y) = 0$ په شان عمل
وکړي چې~$x$ په~$K$ کښې وي؛ خو په همدغو حالاتو کښې مو د هرې بلې
جزوي محاسبه کېدونکې تابع په شان هم عمل پرې کولے شو۔ دا مشاهده موږ ته
اجازه راکوي يوه لا عمومي قضيه بيان کړو چې، په لنډ ډول، وايي د محاسبه
کېدونکو تابعو هېڅ غيرابتذالي خاصيت د پرېکړې وړ نۀ دے۔

په ياد ولرئ چې~$\cfind{0}$، $\cfind{1}$، $\cfind{2}$،~\dots زموږ د
جزوي محاسبه کېدونکو تابعو معياري شمېرنه ده۔

\begin{thm}[د رايس قضيه]
  $C$ د جزوي محاسبه کېدونکو تابعو هر يو سټ وي، او $A =
  \Setabs{n}{\cfind{n} \in C}$ وي۔ که~$A$ محاسبه کېدونکے وي، نو يا
  خو~$C$ تش دے، يا~$C$ د ټولو جزوي محاسبه کېدونکو تابعو سټ دے۔
\end{thm}

\emph{شاخصي سټ} داسې سټ~$A$ دے چې که~$n$ او~$m$ هغه شاخصونه وي چې
يوه تابع ``محاسبه کوي''، نو يا~$n$ او~$m$ دواړه په~$A$ کښې دي، يا
يو هم نۀ دے۔ دا ليدل ستونزمن نۀ دي چې په قضيه کښې سټ~$A$ دا خاصيت
لري۔ برعکس، که~$A$ شاخصي سټ وي او~$C$ د هغو تابعو سټ وي چې دا
شاخصونه ئې محاسبه کوي، نو $A = \Setabs{n}{\cfind{n} \in C}$۔

\begin{explain}
په دې اصطلاح کښې د رايس قضيه له دې وينا سره معادله ده چې هېڅ
غيرابتذالي شاخصي سټ د پرېکړې وړ نۀ دے۔ د قضيې د مطلب د پوهېدو دپاره
ګټوره ده چې د \emph{پروګرامونو} (مثلاً ستاسو په خوښه پروګرامي ژبه
کښې) او د هغو د محاسبه کوونکو تابعو تر منځ توپير روښانه کړو۔ د
پروګرامونو (شاخصونو) په اړه، چې نحوي څيزونه دي، يقيناً داسې پوښتنې
شته چې محاسبه کېدونکې دي: ايا دا پروګرام تر~150 زياتې نښې لري؟ ايا
تر~22 زياتې کرښې لري؟ ايا د ``وايل'' جمله لري؟ ايا رشته ``هېلو
ورلډ'' کله هم د ``پرنټ'' جملې د دليل په توګه راځي؟ د رايس قضيه وايي
چې د پروګرام د \emph{چلند} په اړه هېڅ غيرابتذالي پوښتنه محاسبه
کېدونکې نۀ ده۔ په دې کښې دا پوښتنې هم راځي: ايا پروګرام پر ننوت~$0$
درېږي؟ ايا کله هم درېږي؟ ايا کله هم جفت عدد راګرځوي؟
\end{explain}

\begin{proof}[د رايس د قضيې ثبوت]
فرض کړئ~$C$ نۀ تش دے او نۀ د ټولو جزوي محاسبه کېدونکو تابعو سټ، او
~$A$ په~$C$ کښې د تابعو د شاخصونو سټ وي۔ وبه ښيو چې که~$A$ محاسبه
کېدونکے وای، د درېدنې مسئله مو حل کولے شوه؛ نو~$A$ محاسبه کېدونکے
نۀ دے۔

بې له دې چې عموميت له لاسه ورکړو، فرض کولے شو چې هغه تابع~$f$ چې
هېچرته تعريف شوې نۀ ده، په~$C$ کښې نۀ ده (که نه، په لاندې استدلال
کښې~$C$ او د هغې متمم سره بدل کړئ)۔ $g$ په~$C$ کښې هره تابع وي۔ نظر
دا دے چې که مو~$A$ پرېکړه کولے شو، د هغو شاخصونو چې~$f$ محاسبه کوي
او هغو چې~$g$ محاسبه کوي، توپير مو کولے شو؛ بيا مو همدغه وړتيا د
درېدنې د مسئلې د حل دپاره کارولے شوه۔

طريقه ئې دا ده۔ د نړيوالو جزوي محاسبه کېدونکو تابعو په کارولو دا
تابع تعريفولے شو:
\[
h(x,y) \simeq
\begin{cases}
\text{تعريف شوې نۀ ده} & \text{که $\cfind{x}(x) \fundefined$ وي} \\
g(y) & \text{که داسې نۀ وي۔}
\end{cases}
\]
د~$h$ د محاسبې دپاره لومړے د~$\cfind{x}(x)$ د محاسبې هڅه کوو؛ که
دا محاسبه ودرېږي، بيا~$g(y)$ محاسبه کوو؛ او که \emph{هغه} محاسبه
ودرېږي، وت ئې راګرځوو۔ په لا صوري ډول ليکلے شو:
\[
h(x,y) \simeq \Proj{2}{0}(g(y),\fn{Un}(x,x)).
\]
چېرته چې~$\Proj{2}{0}(z_0, z_1) = z_0$ هغه~$2$-ځاييزه پروجکشن تابع
ده چې~$0$-م دليل راګرځوي، او محاسبه کېدونکې ده۔

بيا~$h$ د جزوي محاسبه کېدونکو تابعو ترکيب دے، او ښے اړخ کټ مټ هغه
وخت تعريف شوے او له~$g(y)$ سره برابر دے چې~$\fn{Un}(x,x)$ او~$g(y)$
دواړه تعريف شوي وي۔

پام وکړئ چې د هر ثابت~$x$ دپاره، که~$\cfind{x}(x)$ تعريف شوې نۀ
وي، نو~$h(x,y)$ د هر~$y$ دپاره تعريف شوې نۀ ده؛ او که~$\cfind{x}(x)$
تعريف شوې وي، نو $h(x,y) \simeq g(y)$۔ نو د~$x$ د هر ثابت قيمت
دپاره~$h(x,y)$ يا کټ مټ د~$f$، يا کټ مټ د~$g$ په شان عمل کوي، او
دا پرېکړه چې~$\cfind{x}(x)$ تعريف شوې ده که نه، د دې پرېکړې سره
برابره ده چې له دې دوو څخه کوم حالت دے۔ خو دا بيا د دې پرېکړې سره
برابره ده چې $h_x(y) \simeq h(x,y)$ په~$C$ کښې ده که نه، او که~$A$
محاسبه کېدونکے وای، همدا کار مو کولے شو۔

په لا صوري ډول، څرنګه چې~$h$ جزوي محاسبه کېدونکې ده، د کوم شاخص~$e$
دپاره له تابع~$\cfind{e}$ سره برابره ده۔ د~$s$-$m$-$n$ قضيې له مخې
داسې بنسټيزه بازګشتي تابع~$s$ شته چې د هر~$x$ دپاره
$\cfind{s(e,x)}(y) \simeq h_x(y)$ وي۔
\textbf{د ژباړې د نسخې سمون (OLCMP-043):} سرچينې دلته د ممکنه
ناتعريفو جزوي تابعو د ارزښتونو تر منځ د عادي برابرۍ نښه کارولې وه؛
د پاراميټري کولو د قضيې له خپل بيان سره سم دلته د جزوي برابرۍ نښه
ساتل شوې ده۔
اوس د هر~$x$ دپاره لرو چې که
$\cfind{x}(x) \fdefined$ وي، نو~$\cfind{s(e,x)}$ هماغه تابع ده چې
~$g$ ده، او ځکه~$s(e,x)$ په~$A$ کښې دے۔ له بلې خوا، که
$\cfind{x}(x) \uparrow$ وي، نو~$\cfind{s(e,x)}$ هماغه تابع ده چې
~$f$ ده، او ځکه~$s(e,x)$ په~$A$ کښې نۀ دے۔ په بله وينا، د هر~$x$
دپاره $x \in K$ هله او يوازې هله چې $s(e,x) \in A$ وي۔ که~$A$
محاسبه کېدونکے وای، ~$K$ هم محاسبه کېدونکے و، چې تناقض دے۔ نو~$A$
محاسبه کېدونکے نۀ دے۔
\end{proof}

د رايس قضيه ډېره پياوړې ده۔ لاندې سمدستي نتيجه ئې د کارونې څو
بېلګې ښيي۔
\begin{cor}
لاندې سټونه د پرېکړې وړ نۀ دي۔
\begin{enumerate}
\item $\Setabs{x}{\text{$17$ د $\cfind{x}$ په قيمتونو کښې دے}}$
\item $\Setabs{x}{\text{$\cfind{x}$ ثابته ده}}$
\item $\Setabs{x}{\text{$\cfind{x}$ هرځاے تعريف شوې ده}}$
\item $\Setabs{x}{\text{هر کله چې $y < y'$، $\cfind{x}(y) \fdefined$، او
    که $\cfind{x}(y') \fdefined$، نو $\cfind{x}(y) < \cfind{x}(y')$}}$
\end{enumerate}
\end{cor}

\begin{proof}
  دا ټول غيرابتذالي شاخصي سټونه دي۔
\end{proof}

\end{document}
